\documentclass[10pt,a4paper]{amsart}
\usepackage[margin=19mm]{geometry}
\usepackage{amsmath,amssymb,mathtools}
\usepackage[scr=rsfs,cal=euler]{mathalfa}
\usepackage{xfrac}
\usepackage[unicode,hidelinks,bookmarks=false]{hyperref}
\newcommand{\ie}{i.e.\ }
\newcommand{\cf}{cf.\ }
\newcommand{\resp}{respectively\ }
\newcommand{\Subsection}{\subsection}
\providecommand{\nobreakdash}{\nobreak\hbox{-}}
\setlength{\emergencystretch}{2em}
\multlinegap0pt
\usepackage{siunitx}
\usepackage{siunitx}
\usepackage{siunitx}
\usepackage{siunitx}
\usepackage{bm}
\usepackage{amssymb}
\usepackage{xcolor}
\usepackage{csquotes}
\usepackage{siunitx}
\newtheorem{assumption}{Assumption}
\newtheorem{lemma}{Lemma}
\newcommand{\bfTheta}{\boldsymbol{\Theta}}
\newcommand{\bfalpha}{\boldsymbol{\alpha}}
\newcommand{\bftheta}{\boldsymbol{\theta}}
\newcommand{\bfxi}{\boldsymbol{\xi}}
\title{On the Asymptotics of Item Selection in Multidimensional Computerized Adaptive Testing}
\author{Seungwon Lee, Xiaoou Li}
\date{Source: arXiv:2608.00201v1 (CC BY 4.0)}
\begin{document}
\maketitle

\noindent\emph{Source and licence.} On the Asymptotics of Item Selection in Multidimensional Computerized Adaptive Testing. Version arXiv:2608.00201v1 (CC BY 4.0), distributed by arXiv under the Creative Commons Attribution 4.0 International licence, http://creativecommons.org/licenses/by/4.0/, as recorded in the arXiv licence metadata for this exact version and shown on the abstract page https://arxiv.org/abs/2608.00201v1. Full licence notice: \texttt{licenses/arxiv-2608.00201-CC-BY-4.0.txt}.\par\smallskip

\noindent\textit{Editorial scope.} Counted unit: the article's lemma lemma:lemma3 (source0.tex lines 1014--1056). The article's complete original proof of it is delivered with it (source0.tex lines 1057--1167, one \texttt{proof} environment of 111 lines). No mathematics is written, repaired or re-derived: the statement and the proof are the article's own lines, in the article's own notation, and no retained line is substituted or re-flowed.  Retained uncounted as the source-local context the proof needs: lines 109--112; lines 227--238; lines 244--251; lines 257--259; lines 269--270; lines 275--276; lines 760--762.  Nothing was omitted from any retained span, so the package carries no omission record.  No retained line contains a citation, so the wrapper carries no bibliography and no bibliography entry.  No retained line contains a figure environment, an image-inclusion command or drawing code, so no image is delivered and the package declares no asset dependency.  Editorial formatting only, in addition: an AMS \texttt{amsart} preamble that restates the article's own declarations and macros used by the retained lines, namely the environments \texttt{assumption}, \texttt{lemma} on separate counters, together with the standard packages those lines need.  The retained environments print in this document as Assumption 1, Lemma 1 in order of appearance, whereas the article prints the same items with the numbering of its own full document.  The complete native source of the article as distributed in the arXiv e-print for this exact version is bundled unchanged as \texttt{sources/206523/source0.tex}.
\begin{equation} \label{e:cumulative_fisher}
\Lambda_{t-1}( \bftheta )=\sum_{s=1}^{t-1} \mathcal{I}_{j_s}(\bftheta), 
\qquad \mathcal{I}_j(\bftheta)=p_j(\bftheta)(1-p_j(\bftheta))\bfalpha_j\bfalpha_j^\top,
\end{equation}

\begin{assumption}[R1: Item-type reuse regime] \label{ass:r1}
Suppose there are $M \in \mathbb{Z}_+$ item types, where $\mathbb{Z}_+$ is the set of positive integers.  
Under R1, an item pool with $M$ calibrated types and $J_r$ operational items is represented as a multiset (i.e., a set that may contain repeating elements)
\begin{equation} \label{e:r1}
\mathcal{J}_{r} 
= \{\,\underbrace{1,\ldots,1}_{J_r/M\text{ times}},\
     \underbrace{2,\ldots,2}_{J_r/M\text{ times}},\ \ldots,\
     \underbrace{M,\ldots,M}_{J_r/M\text{ times}}\,\}, 
\end{equation}
where $J_r$ is a positive integer divisible by $M$ such that $J_r/M\geq T_r$ and $\lim_{r\to\infty}J_r=\infty$.
Letting $r \to \infty$ gives the limiting item pool $\mathcal{J}_{\infty} = \{1, 1, \cdots, 2, 2, \cdots, M, M, \cdots\}$. 
\end{assumption}

\begin{assumption}[R2: Unique-item-parameter regime]
\label{ass:r2}
Suppose the pool contains $J_r \in \mathbb{N}$ calibrated items, where $J_r$ depends on the auxiliary index $r$, satisfies $\lim_{r \to \infty} J_r = \infty$, and is nondecreasing in $r$. For the feasibility of the test, we further assume $J_r \geq T_r$. Then, the item pool for the test of length $T_r$ is given by
\begin{equation} \label{e:r2}
\mathcal{J}_{r} = \{1, \cdots, J_r\},
\end{equation}
Letting $r \to \infty$ gives the limiting item pool $\mathcal{J}_{\infty} = \{1,2,3, \cdots \}$.
\end{assumption}

\begin{assumption} \label{ass:compact} 
The parameter space $\boldsymbol{\Theta}$ is a non-empty compact and convex subset of $\mathbb{R}^d$. The true ability $\bftheta^{\ast}$ is an interior point of $\boldsymbol{\Theta}$.
\end{assumption}

\begin{assumption} \label{ass:information1} (Eigen growth) There exist $c>0$ and $r_0 > 0$ such that $\underline\lambda_{T_r} \geq c$ almost surely for all $r \ge r_0$. Here, $\underline\lambda_{T_r} = \lambda_{\min}(\frac{1}{T_r} \Lambda_{T_r}( \hat{\bftheta}_{T_{r}}^{\mathrm{ML}} ))$, where $\lambda_{\min}(\Sigma)$ denotes the smallest eigenvalue of the matrix $\Sigma$, and $\Lambda_{T_r}(  \hat{\bftheta}_{T_r}^{\mathrm{ML}}  )$ is defined in \eqref{e:cumulative_fisher}. 
\end{assumption}

\begin{assumption} \label{ass:information3} (Uniform boundedness) For the item pool $\mathcal{J}_{\infty}$ defined under either Assumption \ref{ass:r1} or \ref{ass:r2}, $\sup_{j \in \mathcal{J}_{\infty}} \| \bfalpha_j \|_2 < \infty$ and $\sup_{j \in \mathcal{J}_{\infty}} |b_j| < \infty$. 
\end{assumption}

\begin{equation} \label{e:density}
f_{\bftheta, j}(y^j) = \Big( \frac{\exp( \bfalpha_j^T \bftheta - b_j)}{1 + \exp( \bfalpha_j^T \bftheta - b_j)} \Big)^{y^j} \, \Big( \frac{1}{1 + \exp( \bfalpha_j^T \bftheta - b_j)} \Big)^{1 - y^j}, \quad y^j \in \{0,1\}.
\end{equation} 

\begin{lemma} \label{lemma:lemma3}
Suppose Assumptions \ref{ass:r2}, \ref{ass:compact}, \ref{ass:information1}, and \ref{ass:information3} hold. Then the following statements are true.
\begin{enumerate}
\item{There exists a constant $C_{\ell}$ such that 
\begin{equation*}
\sup_{j \in \mathcal{J}_{\infty}, \bftheta \in \bfTheta} | \log f_{\bftheta, j}(y^j) | \leq C_{\ell}.
\end{equation*}}
\item{There exists some constant $C_{\bfalpha} > 0$ from Assumption \ref{ass:information3}, such that
\begin{equation*}
\sup_{j \in \mathcal{J}_{\infty}, \bftheta \in \bfTheta} \| \nabla_{\bftheta} \log f_{\bftheta, j}(y^j) \|_2 \leq C_{\bfalpha}.
\end{equation*}
}
\item{For any $\bftheta_1, \bftheta_2 \in \bfTheta$, there exist some constant $C_{\alpha} > 0$ such that 
\begin{equation*}
| \log f_{\bftheta_1, j}(y^j) - \log f_{\bftheta_2, j}(y^j) | \leq C_{\alpha} \| \bftheta_1 - \bftheta_2 \|_2,
\end{equation*}
where $C_{\alpha}$ is the same constant defined in Assumption \ref{ass:information3}, and does not depend on $j \in \mathcal{J}_{\infty}$.}
\item{For any $\bftheta_1, \bftheta_2 \in \bfTheta$, there exists some constant $C_{\Psi_1} > 0$ such that 
\begin{equation*}
\| \nabla_{\bftheta} \log f_{\bftheta_1, j}(y^j) - \nabla_{\bftheta} \log f_{\bftheta_2, j}(y^j) \|_{2} \leq C_{\Psi_1} \| \bftheta_1 - \bftheta_2 \|_2,
\end{equation*}
where $C_{\Psi_1}$ does not depend on $j \in \mathcal{J}_{\infty}$.}
\item{For any $\bftheta_1, \bftheta_2 \in \bfTheta$, there exists some constant $C_{\Psi_2} > 0$ such that 
\begin{equation*}
\| \nabla_{\bftheta}^2 \log f_{\bftheta_1, j}(y^j) - \nabla_{\bftheta}^2 \log f_{\bftheta_2, j}(y^j) \|_{op} \leq C_{\Psi_2} \| \bftheta_1 - \bftheta_2 \|_2,
\end{equation*}
where $C_{\Psi_2}$ does not depend on $j \in \mathcal{J}_{\infty}$.}
\item{There exists some positive constants $0 < m_1 \le m_2 < \infty$ such that $m_1 \leq | \nabla_{\bfxi_j}^2 \log h_{\bfxi_j, j}(y_j)| \leq m_2$ for all $j \in\mathcal{J}_{\infty}$ and $\bftheta \in \bfTheta$.}
\item{There exists some constant $C_{op}$ such that 
\begin{equation*}
\sup_{j \in \mathcal{J}_{\infty}, \bftheta \in \bfTheta} \| \nabla_{\bftheta}^2 \log f_{\bftheta, j}(y_j) \|_{op} \le C_{op},
\end{equation*}
where $\| \cdot \|_{op}$ is the operator norm.}
\item{For any $r \geq 1$, and any $\bftheta \in \bfTheta$,
\begin{equation*}
\frac{m_1}{m_2} \cdot \frac{1}{T_r} \Lambda_{T_r}(  \hat{\bftheta}_{T_r}^{\mathrm{ML}} ) \preceq 
\frac{1}{T_r} \Lambda_{T_r}(  \bftheta ) 
\preceq
\frac{m_2}{m_1} \cdot \frac{1}{T_r} \Lambda_{T_r}( \hat{\bftheta}_{T_r}^{\mathrm{ML}}  ),
\end{equation*}
for $m_1, m_2$ defined in the previous statement and $\Lambda_{T_r}( \bftheta )$ defined in \eqref{e:cumulative_fisher}.}
\end{enumerate}
\end{lemma}

\begin{proof}
1. From $f_{\bftheta, j}(y^j)$ defined in \eqref{e:density}, $\log f_{\bftheta, j}(y^j)$ is given as
\begin{equation} \label{e:logdensity}
\log f_{\bftheta, j}(y^j) = y^j (\bfalpha_j^T \bftheta - b_j) - \log \Big( 1 + \exp( \bfalpha_j^T \bftheta - b_j )\Big).
\end{equation}
Under Assumptions \ref{ass:compact} and \ref{ass:information3}, we obtain 
\begin{equation*}
\sup_{j \in \mathcal{J}_{\infty}, \bftheta \in \bfTheta} |\bfalpha_j^T \bftheta - b_j| \le C_{\eta},
\end{equation*}
for some $C_{\eta}$. Applying this, we can check
\begin{equation*}
\sup_{j \in \mathcal{J}_{\infty}, \bftheta \in \bfTheta} | \log f_{\bftheta, j}(y^j) | \leq C_{\eta} + \log( 1 + \exp( C_\eta) ).
\end{equation*}
Letting $C_{\ell} = C_{\eta} + \log( 1 + \exp( C_\eta) )$, we prove the statement.

2. Following the log density defined in \eqref{e:logdensity}, taking derivative with respect to $\bftheta$, we obtain
\begin{equation} \label{e:logdenstiy_deriv}
\nabla_{\bftheta} \log f_{\bftheta,j}(y^j) = \Big( y^j - \frac{\exp(\bfalpha_j^T \bftheta - b_j)}{1 + \exp(\bfalpha_j^T \bftheta - b_j)} \Big) \bfalpha_j.
\end{equation}
Note that $\sup_{j \in \mathcal{J}_{\infty}, \bftheta \in \bfTheta} \Big| y^j - \frac{\exp(\bfalpha_j^T \bftheta - b_j)}{1 + \exp(\bfalpha_j^T \bftheta - b_j)} \Big| \le 1$, and $\sup_{j \in \mathcal{J}_{\infty}} \|\bfalpha_j \|_2 \leq C_{\alpha}$ from Assumption \ref{ass:information3}. Thus, 
\begin{equation*}
\sup_{j \in \mathcal{J}_{\infty}, \bftheta \in \bfTheta} \| \nabla_{\bftheta} \log f_{\bftheta, j}(y^j) \|_2 \le C_{\alpha}
\end{equation*}
holds and we prove the statement.

3. Note that $\log f_{\bftheta, j}(y^j)$ is continuous and differentiable function in $\bftheta$. Applying Mean Value Theorem and Cauchy-Schwartz in equality, we have for any $\bftheta_1, \bftheta_2 \in \bfTheta$,
\begin{equation*}
| \log f_{\bftheta_1, j}(y^j) - \log f_{\bftheta_2, j}(y^j) | \le \| \nabla_{\bftheta} \log f_{\tilde{\bftheta}, j}(y^j) \|_2 \| \bftheta_1 - \bftheta_2 \|_2,
\end{equation*}
for some $\tilde{\bftheta}$ between $\bftheta_1$ and $\bftheta_2$. From the previous statement we proved, we obtain $\| \nabla_{\bftheta} \log f_{\tilde{\bftheta}, j}(y^j) \|_2 \le C_{\alpha}$, proving the statement.

4. Let $g_j(\bftheta) = \frac{\exp(\bfalpha_j^T \bftheta - b_j)}{1 + \exp(\bfalpha_j^T \bftheta - b_j)}$. 
Then, $g_j(\bftheta)$ is continuous and differentiable with respect to $\bftheta$. Applying Mean Value Theorem with Assumptions \ref{ass:compact} and \ref{ass:information3}, $g_j(\bftheta)$ is Lipschitz function of $\bftheta$, with some constant $C_1$ not depending on $j \in \mathcal{J}_{\infty}$. For any $\bftheta_1, \bftheta_2 \in \bfTheta$, we can check
\begin{align*}
\| \nabla_{\bftheta} \log f_{\bftheta_1, j}(y^j) - \nabla_{\bftheta} \log f_{\bftheta_2, j}(y^j) \|_2 
&\leq C_1 \|\bftheta_1 - \bftheta_2 \|_2 \| \bfalpha_j \|_2
\\ &\le C_1 \sup_{j \in \mathcal{J}_{\infty}} \| \bfalpha_j \|_2 \|\bftheta_1 - \bftheta_2 \|_2.
\end{align*}
From Assumption \ref{ass:information3}, by letting $C_{\Psi_1} = C_1 C_{\alpha}$, we prove the statement.


5. The first derivative obtained in \eqref{e:logdenstiy_deriv}, is continuous and differentiable with respect to $\bftheta$. Then, the second derivative of $\log f_{\bftheta, j}(y^j)$ is given as
\begin{equation} \label{e:logdensity_deriv2}
\nabla_{\bftheta}^2 \log f_{\bftheta, j}(y^j) = - B_j( \bfalpha_j^T \bftheta) \bfalpha_j \bfalpha_j^T,
\end{equation}
where $B_j(\bfalpha_j^T \bftheta) = \frac{\exp(\bfalpha_j^T \bftheta - b_j)}{\bigl(1+\exp(\bfalpha_j^T \bftheta - b_j) \bigr)^2}$. From $B_j(\bfalpha_j^T \bftheta)$ is continuous and differentiable with respect to $\bftheta$, applying Mean Value Theorem combined with Assumptions \ref{ass:compact} and \ref{ass:information3}, $B_j(\bfalpha_j^T \bftheta)$ is Lipschitz function of $\bftheta$, with some constant $C_1$ not depending on $j \in \mathcal{J}_{\infty}$. Then, for any $\bftheta_1, \bftheta_2 \in \bfTheta$, we obtain
\begin{align*}
\| \nabla_{\bftheta}^2 \log f_{\bftheta_1, j}(y^j) - \nabla_{\bftheta}^2 \log f_{\bftheta_2, j}(y^j) \|_{op} &\leq C_1 \| \bftheta_1 - \bftheta_2 \|_2 \| \bfalpha_j \bfalpha_j^T \|_{op}
\\ &\le C_1 \sup_{j \in \mathcal{J}_{\infty}} \| \bfalpha_j \|_2^2 \| \bftheta_1 - \bftheta_2 \|_2.
\end{align*}
From Assumption \ref{ass:information3}, by defining $C_{\Psi_2} = C_1 C_{\alpha}^2$, we prove the statement.

6. Following $\bfxi_j = \bfalpha_j^T \bftheta$, $\log h_{\bfxi_j, j}(y^j)$ is defined as we can check 
\begin{equation} \label{e:density_reparam}
\log h_{\bfxi_j^{\ast}, j}(y^j) = y^j (\bfxi_j^{\ast} - b_j) - \log( 1 + \exp( \bfxi_j^{\ast} - b_j)).
\end{equation}
Then, the second derivative with respect to $\bfxi_j$ is 
\begin{equation*}
| \nabla_{\bfxi_j}^2 \log h_{\bfxi_j, j}(y_j)| = B_{j}(\bfxi_j).
\end{equation*}
By Assumption \ref{ass:information3} and the compactness of $\bfTheta$, there exist
constants $0 < m_1 \le m_2 < \infty$ such that
\begin{equation*}
m_1 \le B_j(\bfxi_j \bftheta) \le m_2,
\quad
\forall j\in\mathcal J_\infty,\ \forall \bftheta\in\bfTheta.
\end{equation*}
This proves the statement.

7. From \eqref{e:logdensity_deriv2}, we need to verify $\sup_{j \in \mathcal{J}_{\infty}, \bftheta \in \bfTheta} \| B_j(\bfalpha_j^T \bftheta) \bfalpha_j \bfalpha_j^T \|_{op} \leq C_{op}$. Since $\nabla_{\bftheta}^2 \log f_{\bftheta, j}(y_j)$ is rank one matrix, $\| B_j(\bfalpha_j^T \bftheta) \bfalpha_j \bfalpha_j^T \|_{op} = \operatorname{tr}(B_j(\bfalpha_j^T \bftheta) \bfalpha_j \bfalpha_j^T)$. From the $\sup_{j \in \mathcal{J}_{\infty}, \bftheta \in \bfTheta} B_j( \bfxi_j \bftheta)\le m_2$ and Assumption \ref{ass:information3} we obtain
\begin{equation*}
\sup_{j \in \mathcal{J}_{\infty}, \bftheta \in \bfTheta} \operatorname{tr}(B_j(\bfalpha_j^T \bftheta) \bfalpha_j \bfalpha_j^T) \le m_2 C_{\alpha}^2.
\end{equation*}
By defining $C_{op} = m_2 C_{\alpha}^2$, we prove the statement.


8. Recall that
$\Lambda_{T_r}\bigl(\bftheta\bigr)
 = \sum_{s=1}^{T_r}\mathcal I_{j_s}(\bftheta)$
with
$\mathcal I_j(\bftheta) = B_j(\bfalpha_j^T \bftheta)\,
\bfalpha_j\bfalpha_j^T$.
Then
\begin{equation*}
\frac{1}{T_r}\Lambda_{T_r}\bigl(\bftheta \bigr)
= \frac{1}{T_r}\sum_{s=1}^{T_r}
\mathcal I_{j_s}(\hat{\bftheta}^{\mathrm{ML}}_{T_r})
\cdot
\frac{B_{j_s}(\bfalpha_{j_s}^T \bftheta)}{B_{j_s}(\bfalpha_{j_s}^T \hat{\bftheta}^{\mathrm{ML}}_{T_r})}.
\end{equation*}
By the uniform bounds on $B_j$ we have
\begin{equation*}
\frac{m_1}{m_2}
\le
\frac{B_{j_s}(\bfalpha_{j_s}^T \bftheta)}
{B_{j_s}(\bfalpha_{j_s}^T \hat{\bftheta}^{\mathrm{ML}}_{T_r})}
\le
\frac{m_2}{m_1},
\quad s=1,\dots,T_r.
\end{equation*}
Since $\mathcal I_{j_s}(\hat{\bftheta}^{\mathrm{ML}}_{T_r})$ is positive semidefinite,
summing over $s$ yields
\begin{equation*}
\frac{m_1}{m_2}\,\frac{1}{T_r}\Lambda_{T_r}\bigl(\hat{\bftheta}^{\mathrm{ML}}_{T_r}\bigr)
\preceq
\frac{1}{T_r}\Lambda_{T_r}\bigl(\bftheta\bigr)
\preceq
\frac{m_2}{m_1}\,\frac{1}{T_r}\Lambda_{T_r}\bigl( \hat{\bftheta}^{\mathrm{ML}}_{T_r} \bigr).
\end{equation*}
This completes the proof.
\end{proof}

\end{document}
